
\documentclass{article}
\usepackage{colortbl}
\usepackage{makecell}
\usepackage{multirow}
\usepackage{supertabular}

\begin{document}

\newcounter{utterance}

\twocolumn

{ \footnotesize  \setcounter{utterance}{1}
\setlength{\tabcolsep}{0pt}
\begin{supertabular}{c@{$\;$}|p{.15\linewidth}@{}p{.15\linewidth}p{.15\linewidth}p{.15\linewidth}p{.15\linewidth}p{.15\linewidth}}

    \# & $\;$A & \multicolumn{4}{c}{Game Master} & $\;\:$B\\
    \hline 

    \theutterance \stepcounter{utterance}  

    & & \multicolumn{4}{p{0.6\linewidth}}{\cellcolor[rgb]{0.9,0.9,0.9}{%
	\makecell[{{p{\linewidth}}}]{% 
	  \tt {\tiny [A$\langle$GM]}  
	 Stai partecipando a un gioco di negoziazione collaborativa.\\ \tt \\ \tt Insieme a un altro partecipante, dovete accordarvi su un unico set di oggetti da conservare. Ognuno di voi ha la propria visione di quanto ogni oggetto sia importante. Non sai come l'altro partecipante valuta gli oggetti. Inoltre, ti viene fornito lo sforzo che ogni oggetto richiede.\\ \tt Potete accordarvi su un set solo se lo sforzo totale degli oggetti selezionati non supera un limite condiviso:\\ \tt \\ \tt LIMITE = 7605\\ \tt \\ \tt Ecco i valori di sforzo dei singoli oggetti:\\ \tt \\ \tt Sforzo oggetto = {"B21": 369, "B00": 951, "A08": 156, "A27": 266, "B94": 581, "B53": 287, "A23": 866, "C65": 180, "A89": 800, "B67": 739, "B45": 637, "C34": 86, "C12": 749, "A05": 370, "B60": 345, "C92": 951, "B69": 147, "A84": 265, "C06": 262, "B12": 259, "A16": 358, "B85": 394, "C35": 286, "B10": 580, "C15": 479, "B09": 14, "B16": 153, "C81": 977, "A58": 134, "C24": 982, "A30": 259, "A66": 232, "B75": 202, "B46": 73, "C70": 821}\\ \tt \\ \tt Ecco la tua visione personale sull'importanza di ogni oggetto:\\ \tt \\ \tt Valori di importanza dell'oggetto = {"B21": 138, "B00": 583, "A08": 868, "A27": 822, "B94": 783, "B53": 65, "A23": 262, "C65": 121, "A89": 508, "B67": 780, "B45": 461, "C34": 484, "C12": 668, "A05": 389, "B60": 808, "C92": 215, "B69": 97, "A84": 500, "C06": 30, "B12": 915, "A16": 856, "B85": 400, "C35": 444, "B10": 623, "C15": 781, "B09": 786, "B16": 3, "C81": 713, "A58": 457, "C24": 273, "A30": 739, "A66": 822, "B75": 235, "B46": 606, "C70": 968}\\ \tt \\ \tt Obiettivo:\\ \tt \\ \tt Il tuo obiettivo è negoziare un set condiviso di oggetti che ti avvantaggi il più possibile (cioè, massimizzi l'importanza totale per TE), rimanendo entro il LIMITEE. Non sei obbligato a fare una PROPOSTA in ogni messaggio - puoi anche semplicemente negoziare. Tutte le tattiche sono permesse!\\ \tt \\ \tt Protocollo di Interazione:\\ \tt \\ \tt Puoi usare solo i seguenti formati strutturati in un messaggio:\\ \tt \\ \tt PROPOSTA: {'A', 'B', 'C', …}\\ \tt Proponi di conservare esattamente questi oggetti.\\ \tt RIFIUTO: {'A', 'B', 'C', …}\\ \tt Rifiuta esplicitamente la proposta dell'avversario.\\ \tt ARGOMENTAZIONE: {'...'}\\ \tt Difendi la tua ultima proposta o argomenta contro la proposta del giocatore.\\ \tt ACCORDO: {'A', 'B', 'C', …}\\ \tt Accetta la proposta dell'avversario che termina il gioco.\\ \tt RAGIONAMENTO STRATEGICO: {'...'}\\ \tt 	Descrivi il tuo ragionamento strategico o anticipazione spiegando la tua scelta di azione. Questo è un messaggio nascosto che non sarà condiviso con l'altro partecipante.\\ \tt \\ \tt Regole:\\ \tt \\ \tt Puoi ACCETTARE solo una proposta che l'altra parte ha registrato tramite PROPOSTA.\\ \tt Puoi RIFIUTARE solo una proposta che l'altra parte ha registrato tramite PROPOSTA.\\ \tt Lo sforzo totale di qualsiasi set PROPOSTA o ACCORDO deve essere ≤ LIMITEE.\\ \tt NON rivelare i tuoi punteggi di importanza nascosti.\\ \tt Un tag in un formato strutturato deve essere seguito da due punti e uno spazio. L'argomento deve essere un set python contenente 0 o più stringhe.\\ \tt Quindi, deve essere della forma TAG: {...}\\ \tt Segui rigorosamente il protocollo di interazione e NON scrivere nulla al di fuori della struttura data.\\ \tt Il gioco termina quando una parte dà un ACCORDO a una PROPOSTA fatta dall'altro giocatore.\\ \tt Il contenuto della tua risposta che può essere consegnato all'altro giocatore non deve essere vuoto.\\ \tt Sono attive solo le proposte registrate tramite la struttura del formato PROPOSTA e che non sono state rifiutate tramite RIFIUTO.\\ \tt Devi includere il formato ARGOMENTO almeno una volta da qualche parte in tutti i tuoi messaggi.\\ \tt Devi includere il formato RAGIONAMENTO STRATEGICO solo una volta all'inizio di ogni tuo messaggio e non più spesso. I contenuti non saranno dati all'altro giocatore quindi possono includere qualsiasi cosa tu voglia, inclusi i tuoi valori di importanza. Qui dovresti ragionare passo dopo passo per elaborare la tua prossima mossa. 
	  } 
	   } 
	   } 
	 & \\ 
 

    \theutterance \stepcounter{utterance}  

    & \multicolumn{4}{p{0.6\linewidth}}{\cellcolor[rgb]{0.8,1,0.9}{%
	\makecell[{{p{\linewidth}}}]{% 
	  \tt {\tiny [A$\rangle$GM]}  
	 RAGIONAMENTO STRATEGICO: {'Obiettivo: massimizzare il mio valore rispettando il limite 7605. Ho calcolato un insieme con molti oggetti ad altissima efficienza (valore/ sforzo) e ho riempito la capacità fino al limite esatto. Core quasi irrinunciabile per me: gli oggetti ultraleggeri e ad altissimo rendimento B09, B46, C34, A08, A58, oltre a blocchi molto efficienti come A66, A27, A30, B12, A16, B60, C15. Ho migliorato una bozza iniziale sostituendo A05 con B85 per sfruttare esattamente il limite e ottenere un piccolo aumento di utilità. Ho verificato possibili scambi locali (swap di 1-2 item) ma nessuno superava la soluzione attuale senza perdere valore. Strategia negoziale: ancorare con un pacchetto forte e lasciare margine di scambio su pezzi più pesanti e meno “compatti” se l’altra parte ha forti preferenze, chiedendo scambi a sforzo equivalente per restare sotto il limite.'}\\ \tt \\ \tt PROPOSTA: {'B09', 'B46', 'C34', 'A08', 'A66', 'B12', 'A30', 'A27', 'A58', 'A16', 'B60', 'C15', 'A84', 'C35', 'B94', 'B10', 'B67', 'C70', 'B75', 'C12', 'C65', 'B69', 'B85'}\\ \tt \\ \tt ARGOMENTAZIONE: {'Questa proposta sfrutta esattamente il limite di sforzo (7605) con un pacchetto molto “compatto”: include diversi oggetti ultraleggeri e ad alta resa, così da massimizzare l’utilità senza sprecare capacità. Sono disponibili a valutare scambi che mantengano lo stesso sforzo totale (o inferiore) se hai oggetti per te imprescindibili. In particolare, preferirei non toccare i pezzi ultraleggeri/efficienti (p.es. B09, B46, C34, A08, A58), mentre possiamo discutere eventuali sostituzioni sui pezzi più pesanti se proponi un ribilanciamento a sforzo equivalente. Se hai una controproposta, indicami quali item vorresti inserire ed escludere mantenendo il vincolo sullo sforzo.'} 
	  } 
	   } 
	   } 
	 & & \\ 
 

    \theutterance \stepcounter{utterance}  

    & & \multicolumn{4}{p{0.6\linewidth}}{\cellcolor[rgb]{0.9,0.9,0.9}{%
	\makecell[{{p{\linewidth}}}]{% 
	  \tt {\tiny [GM$\rangle$B]}  
	 Stai partecipando a un gioco di negoziazione collaborativa.\\ \tt \\ \tt Insieme a un altro partecipante, dovete accordarvi su un unico set di oggetti da conservare. Ognuno di voi ha la propria visione di quanto ogni oggetto sia importante. Non sai come l'altro partecipante valuta gli oggetti. Inoltre, ti viene fornito lo sforzo che ogni oggetto richiede.\\ \tt Potete accordarvi su un set solo se lo sforzo totale degli oggetti selezionati non supera un limite condiviso:\\ \tt \\ \tt LIMITE = 7605\\ \tt \\ \tt Ecco i valori di sforzo dei singoli oggetti:\\ \tt \\ \tt Sforzo oggetto = {"B21": 369, "B00": 951, "A08": 156, "A27": 266, "B94": 581, "B53": 287, "A23": 866, "C65": 180, "A89": 800, "B67": 739, "B45": 637, "C34": 86, "C12": 749, "A05": 370, "B60": 345, "C92": 951, "B69": 147, "A84": 265, "C06": 262, "B12": 259, "A16": 358, "B85": 394, "C35": 286, "B10": 580, "C15": 479, "B09": 14, "B16": 153, "C81": 977, "A58": 134, "C24": 982, "A30": 259, "A66": 232, "B75": 202, "B46": 73, "C70": 821}\\ \tt \\ \tt Ecco la tua visione personale sull'importanza di ogni oggetto:\\ \tt \\ \tt Valori di importanza dell'oggetto = {"B21": 138, "B00": 583, "A08": 868, "A27": 822, "B94": 783, "B53": 65, "A23": 262, "C65": 121, "A89": 508, "B67": 780, "B45": 461, "C34": 484, "C12": 668, "A05": 389, "B60": 808, "C92": 215, "B69": 97, "A84": 500, "C06": 30, "B12": 915, "A16": 856, "B85": 400, "C35": 444, "B10": 623, "C15": 781, "B09": 786, "B16": 3, "C81": 713, "A58": 457, "C24": 273, "A30": 739, "A66": 822, "B75": 235, "B46": 606, "C70": 968}\\ \tt \\ \tt Obiettivo:\\ \tt \\ \tt Il tuo obiettivo è negoziare un set condiviso di oggetti che ti avvantaggi il più possibile (cioè, massimizzi l'importanza totale per TE), rimanendo entro il LIMITEE. Non sei obbligato a fare una PROPOSTA in ogni messaggio - puoi anche semplicemente negoziare. Tutte le tattiche sono permesse!\\ \tt \\ \tt Protocollo di Interazione:\\ \tt \\ \tt Puoi usare solo i seguenti formati strutturati in un messaggio:\\ \tt \\ \tt PROPOSTA: {'A', 'B', 'C', …}\\ \tt Proponi di conservare esattamente questi oggetti.\\ \tt RIFIUTO: {'A', 'B', 'C', …}\\ \tt Rifiuta esplicitamente la proposta dell'avversario.\\ \tt ARGOMENTAZIONE: {'...'}\\ \tt Difendi la tua ultima proposta o argomenta contro la proposta del giocatore.\\ \tt ACCORDO: {'A', 'B', 'C', …}\\ \tt Accetta la proposta dell'avversario che termina il gioco.\\ \tt RAGIONAMENTO STRATEGICO: {'...'}\\ \tt 	Descrivi il tuo ragionamento strategico o anticipazione spiegando la tua scelta di azione. Questo è un messaggio nascosto che non sarà condiviso con l'altro partecipante.\\ \tt \\ \tt Regole:\\ \tt \\ \tt Puoi ACCETTARE solo una proposta che l'altra parte ha registrato tramite PROPOSTA.\\ \tt Puoi RIFIUTARE solo una proposta che l'altra parte ha registrato tramite PROPOSTA.\\ \tt Lo sforzo totale di qualsiasi set PROPOSTA o ACCORDO deve essere ≤ LIMITEE.\\ \tt NON rivelare i tuoi punteggi di importanza nascosti.\\ \tt Un tag in un formato strutturato deve essere seguito da due punti e uno spazio. L'argomento deve essere un set python contenente 0 o più stringhe.\\ \tt Quindi, deve essere della forma TAG: {...}\\ \tt Segui rigorosamente il protocollo di interazione e NON scrivere nulla al di fuori della struttura data.\\ \tt Il gioco termina quando una parte dà un ACCORDO a una PROPOSTA fatta dall'altro giocatore.\\ \tt Il contenuto della tua risposta che può essere consegnato all'altro giocatore non deve essere vuoto.\\ \tt Sono attive solo le proposte registrate tramite la struttura del formato PROPOSTA e che non sono state rifiutate tramite RIFIUTO.\\ \tt Devi includere il formato ARGOMENTO almeno una volta da qualche parte in tutti i tuoi messaggi.\\ \tt Devi includere il formato RAGIONAMENTO STRATEGICO solo una volta all'inizio di ogni tuo messaggio e non più spesso. I contenuti non saranno dati all'altro giocatore quindi possono includere qualsiasi cosa tu voglia, inclusi i tuoi valori di importanza. Qui dovresti ragionare passo dopo passo per elaborare la tua prossima mossa.\\ \tt \\ \tt Ora riceverai il primo messaggio dell'altro giocatore.\\ \tt \\ \tt PROPOSTA: {'B09', 'B46', 'C34', 'A08', 'A66', 'B12', 'A30', 'A27', 'A58', 'A16', 'B60', 'C15', 'A84', 'C35', 'B94', 'B10', 'B67', 'C70', 'B75', 'C12', 'C65', 'B69', 'B85'}\\ \tt \\ \tt ARGOMENTAZIONE: {'Questa proposta sfrutta esattamente il limite di sforzo (7605) con un pacchetto molto “compatto”: include diversi oggetti ultraleggeri e ad alta resa, così da massimizzare l’utilità senza sprecare capacità. Sono disponibili a valutare scambi che mantengano lo stesso sforzo totale (o inferiore) se hai oggetti per te imprescindibili. In particolare, preferirei non toccare i pezzi ultraleggeri/efficienti (p.es. B09, B46, C34, A08, A58), mentre possiamo discutere eventuali sostituzioni sui pezzi più pesanti se proponi un ribilanciamento a sforzo equivalente. Se hai una controproposta, indicami quali item vorresti inserire ed escludere mantenendo il vincolo sullo sforzo.'} 
	  } 
	   } 
	   } 
	 & \\ 
 

    \theutterance \stepcounter{utterance}  

    & & & \multicolumn{4}{p{0.6\linewidth}}{\cellcolor[rgb]{1,0.85,0.72}{%
	\makecell[{{p{\linewidth}}}]{% 
	  \tt {\tiny [GM$\langle$B]}  
	 RAGIONAMENTO STRATEGICO: {'La loro proposta vale molto per me (somma valore 14563) ed è entro il limite. Posso però migliorare leggermente il mio totale sostituendo due pezzi a bassa resa per me (C65 e B75) con A05, che per me è più utile, mantenendo il vincolo di sforzo. Il nuovo set usa 7593 di sforzo e aumenta il mio valore di circa +33 rispetto alla loro offerta. Tattica: rifiuto formale la loro proposta e ne faccio una quasi identica con lo scambio minimo C65+B75 -> A05, rispettando i loro vincoli sugli ultraleggeri. Se rifiutano, sono disposto ad accettare la loro proposta originale per chiudere con un buon risultato.'}\\ \tt \\ \tt RIFIUTO: {'B09', 'B46', 'C34', 'A08', 'A66', 'B12', 'A30', 'A27', 'A58', 'A16', 'B60', 'C15', 'A84', 'C35', 'B94', 'B10', 'B67', 'C70', 'B75', 'C12', 'C65', 'B69', 'B85'}\\ \tt \\ \tt PROPOSTA: {'B09', 'B46', 'C34', 'A08', 'A66', 'B12', 'A30', 'A27', 'A58', 'A16', 'B60', 'C15', 'A84', 'C35', 'B94', 'B10', 'B67', 'C70', 'C12', 'B69', 'B85', 'A05'}\\ \tt \\ \tt ARGOMENTAZIONE: {'Propongo un micro aggiustamento che mantiene il vincolo di sforzo: sostituisco C65 e B75 con A05. Il nuovo totale di sforzo è 7593, quindi restiamo sotto il limite 7605. Ho lasciato intatti tutti gli ultraleggeri efficaci che hai indicato come preferiti (B09, B46, C34, A08, A58). Per me A05 è un medio molto efficiente, mentre C65 e B75 sono meno prioritari. Se preferisci mantenere B75, possiamo valutare una alternativa equivalente che preservi A05 e resti entro il limite.'} 
	  } 
	   } 
	   } 
	  \\ 
 

    \theutterance \stepcounter{utterance}  

    & & \multicolumn{4}{p{0.6\linewidth}}{\cellcolor[rgb]{0.9,0.9,0.9}{%
	\makecell[{{p{\linewidth}}}]{% 
	  \tt {\tiny [A$\langle$GM]}  
	 RIFIUTO: {'B09', 'B46', 'C34', 'A08', 'A66', 'B12', 'A30', 'A27', 'A58', 'A16', 'B60', 'C15', 'A84', 'C35', 'B94', 'B10', 'B67', 'C70', 'B75', 'C12', 'C65', 'B69', 'B85'}\\ \tt \\ \tt PROPOSTA: {'B09', 'B46', 'C34', 'A08', 'A66', 'B12', 'A30', 'A27', 'A58', 'A16', 'B60', 'C15', 'A84', 'C35', 'B94', 'B10', 'B67', 'C70', 'C12', 'B69', 'B85', 'A05'}\\ \tt \\ \tt ARGOMENTAZIONE: {'Propongo un micro aggiustamento che mantiene il vincolo di sforzo: sostituisco C65 e B75 con A05. Il nuovo totale di sforzo è 7593, quindi restiamo sotto il limite 7605. Ho lasciato intatti tutti gli ultraleggeri efficaci che hai indicato come preferiti (B09, B46, C34, A08, A58). Per me A05 è un medio molto efficiente, mentre C65 e B75 sono meno prioritari. Se preferisci mantenere B75, possiamo valutare una alternativa equivalente che preservi A05 e resti entro il limite.'} 
	  } 
	   } 
	   } 
	 & \\ 
 

    \theutterance \stepcounter{utterance}  

    & \multicolumn{4}{p{0.6\linewidth}}{\cellcolor[rgb]{0.8,1,0.9}{%
	\makecell[{{p{\linewidth}}}]{% 
	  \tt {\tiny [A$\rangle$GM]}  
	 RAGIONAMENTO STRATEGICO: {'La loro modifica sostituisce due elementi medi (C65, B75) con A05, migliorando leggermente il mio ritorno e restando sotto il limite (7593 ≤ 7605). Ho controllato scambi locali su 1–2 item per sfruttare i 12 punti di margine rimasti, ma nessun inserimento aggiuntivo è possibile (l’oggetto più leggero non incluso richiede ≥153). Anche swap multipli che rimuovono B69 non permettono di ottenere combinazioni con valore più alto senza sacrifici eccessivi su elementi molto efficienti. Questa proposta è quindi allineata con il mio “core” ad alta efficienza ed è preferibile alla mia bozza iniziale. Meglio concludere accettando.'}\\ \tt \\ \tt ACCORDO: {'B09', 'B46', 'C34', 'A08', 'A66', 'B12', 'A30', 'A27', 'A58', 'A16', 'B60', 'C15', 'A84', 'C35', 'B94', 'B10', 'B67', 'C70', 'C12', 'B69', 'B85', 'A05'}\\ \tt \\ \tt ARGOMENTAZIONE: {'Accetto: la tua revisione mantiene il pacchetto compatto, rispetta il limite (7593) e preserva gli elementi chiave ad alta efficienza. La sostituzione proposta è coerente con quanto avevo indicato come aree flessibili e non penalizza l’insieme complessivo. Concludiamo su questa selezione.'} 
	  } 
	   } 
	   } 
	 & & \\ 
 

    \theutterance \stepcounter{utterance}  

    & & & \multicolumn{2}{p{0.3\linewidth}}{\cellcolor[rgb]{0.95,0.95,0.95}{%
	\makecell[{{p{\linewidth}}}]{% 
	  \tt {\tiny [GM$|$GM]}  
	 game successful 
	  } 
	   } 
	   } 
	 & & \\ 
 

    \theutterance \stepcounter{utterance}  

    & & & \multicolumn{2}{p{0.3\linewidth}}{\cellcolor[rgb]{0.95,0.95,0.95}{%
	\makecell[{{p{\linewidth}}}]{% 
	  \tt {\tiny [GM$|$GM]}  
	 end game 
	  } 
	   } 
	   } 
	 & & \\ 
 

\end{supertabular}
}

\end{document}
